\documentclass[11pt,a4paper]{article}
\usepackage[T1]{fontenc}
\usepackage{lmodern,amsmath,amssymb,amsthm,mathtools,booktabs,microtype}
\usepackage[margin=27mm]{geometry}
\usepackage[colorlinks=true,linkcolor=blue,citecolor=blue,urlcolor=blue]{hyperref}
\setlength{\emergencystretch}{2em}
\newtheorem{theorem}{Theorem}[section]
\newtheorem{lemma}[theorem]{Lemma}
\newtheorem{proposition}[theorem]{Proposition}
\newtheorem{corollary}[theorem]{Corollary}
\newtheorem{definition}[theorem]{Definition}
\newcommand{\Q}{\mathbb Q}\newcommand{\Z}{\mathbb Z}
\DeclareMathOperator{\Tr}{Tr}\DeclareMathOperator{\Nm}{N}
\title{A guarded rational-square criterion for normalized elliptic detectors}
\author{}\date{16 November 2026}
\begin{document}\maketitle
\begin{abstract}
We give a finite quadratic test excluding negative odd detector poles
when every odd valuation of the physical weight is positive. The test
is strictly positive at every prepared character of the point algebra
and has an exact square value at the algebraic reference product
$W=1$. For the conditional sources, a quadratic detector product with
faster decay gives local square values at every fixed preparation
place. The exact multiplication family also satisfies the local test
at its selected split contacts. Two simultaneous rational square values
on actual guarded detector tuples would give twenty-one total variables
for the general source and nine for either conditional source. That
global realization statement is not proved; the established defining
counts remain twenty-five and thirteen.
\end{abstract}

\section{Detector hypotheses and arithmetic scope}

Let $K=\Q(\sqrt2)$ and let $S$ be the source's finite exceptional
set of rational primes, containing all primes at most $76176$.
Every valuation above
$p\notin S$ is normalized by $v(p)=1$. We use an actual field-valued
detector product $W$, a nonzero rational physical weight $w$, and
\begin{equation}\label{eq:native}
 V=\Tr_{K/\Q}(W^{-1}),\qquad
 J=\Nm_{K/\Q}(W^{-1}),\qquad b=\frac{4w}{V^2}.
\end{equation}
All point, difference, rational-function and trace denominators are
guarded. The following hypotheses concern the guarded rational source.
\begin{enumerate}
\item Every odd valuation of $w$ outside $S$ is positive.
\item At an odd valuation of $w$, each local reciprocal slot of $W$
is integral and reduces to either zero or one.
\item At a prepared specialization, every geometric character of the
finite point algebra is real, $w>0$, and both real slots of $W$ are
at least one.
\end{enumerate}

The normalized general source in~\cite{Normalized} satisfies these
hypotheses. Its physical weight is
\[
             w=\frac{u}{(1+u)^2},\qquad u=\rho R,
\]
with $1+u$ guarded. Its field factors use this physical weight in
every denominator, and its scalar numerators are
\[
 G_k=R^{\kappa_k+1}\Phi_R(S_k),\qquad
 \Phi_R(s)=\frac{Rs}{1+Rs}.
\]
The nine point roles and common slope are those of the general
source~\cite{General,Unit}. All assertions about their integer input
projection retain the arithmetic premises specified there.

For the conditional sources, put
\begin{equation}\label{eq:conditionalweight}
 F(t)=\frac{t(t+1)(2t+1)(3t+1)(4t+1)(5t+2)}{120},\qquad
 \delta=t-1,\qquad w=\frac{F(t)}{4\delta^8}.
\end{equation}
The six actual field tests and the faster quadratic product are
specified in Section~\ref{sec:conditional}. The ordinary and weak
abc assumptions, subgroup input and remaining arithmetic premises
are exactly those of the corresponding conditional sources
\cite{Conditional,Normalized}.

\begin{lemma}\label{lem:weight}
The displayed general weight has positive odd valuations. The same
is true of~\eqref{eq:conditionalweight} at every rational guarded
argument outside the conditional exceptional set.
\end{lemma}
\begin{proof}
If $n=v_p(u)\ne0$, then
\[
 v_p\left(\frac{u}{(1+u)^2}\right)=|n|.
\]
When $n=0$ the valuation is even. For the conditional weight,
$v_p(t)<0$ gives $v_p(w)=-2v_p(t)$, since $F$ is monic of degree
six outside $S$. If $v_p(t-1)>0$, then $F(t)\equiv F(1)=7$ is a
unit and $v_p(w)=-8v_p(t-1)$. Otherwise $t$ and all coefficients
are integral and the denominator is a unit. Every odd valuation is
therefore positive.
\end{proof}

\section{The finite rational-square test}

Set
\begin{equation}\label{eq:criterion}
 g=1+w,\qquad Q=1-\frac Vg+\frac J{g^2},\qquad
 \mathcal R_a=Q\left(1+\frac{a^2(1-V/2)}w\right),
 \qquad a\in\Z[S^{-1}]^\times.
\end{equation}
The coefficient $a$ is fixed, so it is not an additional rational
coordinate. Different fixed coefficients may be used in the two blocks.

\begin{theorem}\label{thm:criterion}
Under the three detector hypotheses, a rational-square value of
$\mathcal R_a$ excludes every negative odd valuation of $b$ outside
$S$. At every prepared point-algebra specialization,
$\mathcal R_a$ is strictly positive at every geometric character and
is a whole-algebra unit. Its value at the algebraic reference product
$W=1$ is
\begin{equation}\label{eq:baseline}
                  \mathcal R_a=\frac{w^2}{(1+w)^2}.
\end{equation}
\end{theorem}
\begin{proof}
Suppose $v_p(b)<0$ is odd. Then $\nu=v_p(w)>0$ is odd and
$v_p(V)>\nu/2>0$. The local reciprocal slots of $W$ reduce to zero
or one. Their trace reduces to zero, so both slots reduce to zero:
at a split prime this uses $p\ne2$, and at an inert prime the trace
of the residue one is two. Hence $J\equiv0$, $g\equiv1$ and
$Q\equiv1\pmod p$. Since $a$ is a local unit, the second factor
of~\eqref{eq:criterion} has valuation $-\nu$. Thus
$v_p(\mathcal R_a)=-\nu$ is odd, contrary to rational squareness.

At a prepared character, write the two reciprocal slots of $W$ as
$u_+,u_-\in(0,1]$. Then $0<V=u_++u_-\le2$, and
\[
 Q=\left(1-\frac{u_+}{g}\right)
   \left(1-\frac{u_-}{g}\right)>0
\]
because $g>1$. The second factor of~\eqref{eq:criterion} is at
least one. Every character is real, so the strict inequalities give
the whole-algebra unit assertion. Finally, $W=1$ gives $V=2$, $J=1$,
and hence~\eqref{eq:baseline}.
\end{proof}

For coefficient construction write $W=A+B\sqrt2$ and $N=A^2-2B^2$.
The exact rational expression is
\begin{equation}\label{eq:cleared}
 \mathcal R_a=
 \frac{(g^2N-2gA+1)\bigl((w+a^2)N-a^2A\bigr)}
 {wg^2N^2}.
\end{equation}
In particular, the square test introduces no denominator divisor beyond
$w$, $1+w$ and the retained product norm. All are nonzero at every
prepared character.

The choice $a=1$ has a useful norm--trace form. Define
\[
 E=1-(gW)^{-1}=e_0+e_1\sqrt2,\qquad
 \omega=\frac{w}{g^2}.
\]
Then $Q=\Nm(E)$, $e_0=1-V/(2g)$, and
\begin{equation}\label{eq:flat}
 \mathcal R_1=
 \frac{e_0\Nm(E)}{\omega g}
 =Q\left(1+\frac{1-V/2}{w}\right).
\end{equation}
Equation~\eqref{eq:flat} always uses the actual detector trace and
norm. A parametrization with freely assigned $e_0,e_1$ would not
establish a realization of this test.

\begin{proposition}\label{prop:adaptive}
The fixed coefficient $a$ in Theorem~\ref{thm:criterion} may be replaced
by a rational function of the retained source coordinates, provided
that its denominators are guarded and $v_p(a)\le0$ whenever
$p\notin S$ and $b$ has negative odd valuation. No new unrestricted
coordinate is needed for such a function.

For example, let $M$ be a fixed nonzero $S$-unit and let $P_0,P_1$
be polynomials in $w,V,J$ over $\Z[S^{-1}]$, whose constant terms
are $S$-units. Then
\[
                    a=M\frac{P_0(w,V,J)}{P_1(w,V,J)}
\]
has the required local unit property on its guarded domain.
If $P_1$ has nonnegative coefficients and a positive constant term,
its denominator is nonzero at every prepared character.
\end{proposition}
\begin{proof}
At a hypothetical bad prime put $\nu=v_p(w)>0$. The proof of
Theorem~\ref{thm:criterion} gives $Q\equiv1$ and
$1-V/2\equiv1\pmod p$. The unique lowest valuation in the second
factor is now
\[
                         2v_p(a)-\nu<0,
\]
which is odd. This proves the first assertion. For the displayed
polynomial family, $(w,V,J)$ reduces to $(0,0,0)$ at a hypothetical
bad prime, so both polynomial values are local units. At a prepared
character $w,V,J>0$, which proves the last assertion. Squaring the
real value of $a$ preserves the positivity argument.
\end{proof}

\begin{corollary}\label{cor:outer}
Let $c$ be an existing guarded rational source function. Then
\[
                              a=1+wc^2
\]
satisfies the hypothesis of Proposition~\ref{prop:adaptive}. Products
of such factors are also admissible. These choices add no unrestricted
coordinate and no additional nonzero guard beyond the denominators
already retained for the chosen source functions. At a prepared real
character each factor is at least one.
\end{corollary}
\begin{proof}
At an odd physical valuation $\nu>0$, if $c\ne0$ the two summands
have valuations zero and the odd integer $\nu+2v_p(c)$. Hence
\[
                    v_p(1+wc^2)=\min\{0,\nu+2v_p(c)\}\le0.
\]
For $c=0$ the factor is one. Product valuations add, and prepared
real positivity follows from $w>0$.
\end{proof}

The rationality of $c$ is essential to this parity calculation. An
unconstrained finite-extension root may have a half-integral valuation.
Introducing a new finite equation for $c$ therefore requires its own
rational-character and completeness argument; it does not provide an
arbitrary extra rational parameter without a coordinate cost.

For $a=1$ the actual arithmetic condition also has a direct elliptic
form. Write $W^{-1}=u+v\sqrt2$, and define the prescribed rational
quantities
\[
              \epsilon=1-u,\qquad x=1+\epsilon/w,\qquad q=v/w.
\]
Then the exact identity
\begin{equation}\label{eq:elliptic}
          \mathcal R_1=\frac{w^2}{(1+w)^2}\,x(x^2-2q^2)
\end{equation}
shows that the square condition is equivalent to
\[
                  y^2=x^3-2q^2x,\qquad
                  y=\frac{1+w}{w}\sqrt{\mathcal R_1}.
\]
For $q\ne0$ this is a rational $j=1728$ twist with its abscissa and
twist coefficient prescribed by the actual detector. The identity
remains valid at $q=0$; interpreting that degenerate equation as an
elliptic curve is neither required nor used. No rational point with
the prescribed coordinates, or at both detector arguments, follows
from the identity alone.

\section{The faster conditional product and its symmetry}
\label{sec:conditional}

Use the source curve, fixed point $P$, and three point roles $A,H,U$
from~\cite{Conditional}. Write $f=1/x$, $z=-x/y$ and
\[
 r=\frac{z(U)}{z(H)},\qquad
 s=\frac{z(U-A)}{z(H-P)}.
\]
Conjugation in $K/\Q$ is denoted by a bar. Define
\begin{equation}\label{eq:tests}
 \begin{gathered}
 C_1=f(H),\quad C_2=f(U),\quad
 C_3=\overline{f(H-P)},\quad C_4=\overline{f(U-A)},\\
 C_5=r-\bar s,\qquad C_6=r+\bar s-2X.
 \end{gathered}
\end{equation}
The product used here is
\begin{equation}\label{eq:fastproduct}
                  W(t)=\prod_{i=1}^6\left(1+\frac{C_i^2}{F(t)}\right).
\end{equation}
The physical weight remains~\eqref{eq:conditionalweight}. The two
arguments are the actual $t_0=\lambda$ and $t_1=a_1\lambda$, with
the source's unchanged input-dependent height scale $a_1$.

\begin{proposition}\label{prop:conditional}
At an odd physical weight valuation $\nu>0$ outside $S$, every local
reciprocal slot of~\eqref{eq:fastproduct} is integral and reduces to
zero or one. If the corresponding physical input has no negative odd
valuation, all six original source tests have positive valuation at
one place above the prime. Thus the original local extraction and
height argument apply.
\end{proposition}
\begin{proof}
Lemma~\ref{lem:weight} shows that $t$ is integral and $\delta$ is a
unit, so $v_p(F)=\nu$. Each factor has valuation
\[
                    \min\{0,2v(C_i)-\nu\}\le0.
\]
The two displayed terms have distinct parity. If the minimum is zero,
the factor reduces to one; otherwise its inverse reduces to zero.

Put $q=\max_{\mathfrak p\mid p}v_{\mathfrak p}(W)\le0$.
Then $v_p(V)\ge-q$. Nonnegative odd valuation of $b$ gives
$1\le v_p(b)\le\nu+2q$, hence $q>-\nu/2$. At a place attaining
$q$, every individual factor valuation is at least $q$, so
\[
                         v(C_i)>\nu/4>0.
\]
The original last test is $(C_5+C_6)/2=r-X$. Since $2$ is a unit,
both original comparisons have positive valuation. All six original
tests therefore meet their order-one thresholds, which are exactly
the hypotheses used by the conditional local extraction proof.
\end{proof}

The ordered source completeness also transfers to this product: the
selected-contact orders are unchanged, while its fixed-place and real
decay are stronger. The full retained-source statement, including the
unchanged thirteen-variable encoding with pole blocks, is given in
\cite{Normalized}.

\begin{proposition}\label{prop:symmetry}
On the common guarded rational-function domain, the involution
\[
                         \tau:(H,U)\longmapsto(P-H,A-U)
\]
satisfies $W\circ\tau=\bar W$. Consequently $V,J,Q$ and
$\mathcal R_a$ descend under $\tau$ for fixed rational $a$.
The exact multiplication family $A=[X]P$, $U=[X]H$ is preserved.
\end{proposition}
\begin{proof}
The involution interchanges $r$ and $s$. It sends the first four
tests to the conjugates of $C_3,C_4,C_1,C_2$, respectively, and
\[
                C_5\circ\tau=-\bar C_5,\qquad
                C_6\circ\tau=\bar C_6.
\]
Their squares give the stated product identity. The trace and norm
are invariant under conjugation. Finally,
$[X](P-H)=A-U$ on the exact multiplication family.
\end{proof}

The statement concerns the common open domain of the charts and their
translates. Quotient points still have to lift to actual source points,
including the square-abscissa lift and every original guard. No
coordinate saving is inferred from the symmetry alone.

\section{Local realization of the actual square radicands}

\begin{lemma}[Stability at the reference product]\label{lem:stability}
Let $p$ be any finite prime. For fixed $w_0\in\Q_p\setminus\{0,-1\}$,
if $w\to w_0$ and both slots of $W\to1$, then
\[
                  \mathcal R_a\longrightarrow
                           \frac{w_0^2}{(1+w_0)^2}.
\]
Thus the actual radicand is a nonzero square in $\Q_p$ throughout a
sufficiently small neighborhood. If instead $w\to0$ and
$W_\pm-1=o(w)$, then
\begin{equation}\label{eq:rapidsquare}
             \frac{\mathcal R_a}{(w/(1+w))^2}\longrightarrow1.
\end{equation}
The same local-square conclusion holds. Here $a$ is fixed, or tends
to a finite $p$-adic limit.
\end{lemma}
\begin{proof}
The first statement is continuity and~\eqref{eq:baseline}. For the
second, write $W_\pm=1+\varepsilon_\pm$. Then
\[
 \frac{1-1/(gW_\pm)}{w/g}
             =1+\frac{\varepsilon_\pm}{w(1+\varepsilon_\pm)}
             \longrightarrow1,
 \qquad 1-V/2=o(w).
\]
Multiplying the two first identities and using the last gives
\eqref{eq:rapidsquare}. Nonzero squares form an open subgroup of
$\Q_p^\times$: a sufficiently close quotient is one modulo $p$ for
odd $p$, and one modulo eight at $p=2$.
\end{proof}

For the faster conditional product, fixed actual tests give, uniformly
on bounded finite-place or compact real neighborhoods,
\begin{equation}\label{eq:infinity}
 w(t)\sim\frac1{4t^2},\qquad
 W(t)-1=O(t^{-6})=O(w(t)^3)\qquad(t\longrightarrow\infty).
\end{equation}
Consequently Lemma~\ref{lem:stability} applies at every fixed finite
preparation place, including $2,3,17$, simultaneously to the two
arguments $t_0,t_1$. Both scales are fixed before the slope is chosen.
For the normalized general source, the finite preparation has
$w_j\to\rho/(1+\rho)^2\ne0,-1$ and $W_j\to1$, so the first
part of that lemma applies instead.

\begin{lemma}[Fresh split contacts]\label{lem:contacts}
Let $p\notin S$ split in $K$. Suppose the actual native preparation
has
\[
 v_p(w)=1,\qquad W_{\rm good}=1+O(p^2),\qquad
                         W_{\rm bad}^{-1}=O(p^2).
\]
If $v_p(a)\le0$, then
\[
                         \mathcal R_a/a^2\equiv1/2\pmod p,
\]
and the actual radicand is a nonzero square in $\Q_p$.
\end{lemma}
\begin{proof}
One has $1-(gW_{\rm good})^{-1}=w+O(p^2)$ and
$1-(gW_{\rm bad})^{-1}=1+O(p^2)$. Hence
$Q=w+O(p^2)$ and $V=1+O(p^2)$. Now
\[
 \frac{\mathcal R_a}{a^2}
        =\frac Q{a^2}+\frac Qw(1-V/2).
\]
The first term is divisible by $p$, giving the stated residue.
Since $p$ splits in
$\Q(\sqrt2)$, the element $2$ is a local square, and Hensel's
lemma gives the conclusion.
\end{proof}

The exact multiplication family in the conditional source supplies
these orders: the first four successful tests have valuation at least
two, the two comparisons have valuation at least four, and $v_p(F)=1$.
Thus the successful product is $1+O(p^3)$. Excluding the fixed
coordinate and denominator support of $P$ and $A$ makes the opposite
reciprocal product $O(p^2)$. The strengthened scalar normalizers in
the general source give the same sufficient orders~\cite{Normalized}.

\begin{lemma}[Conditional pole places]\label{lem:poles}
Suppose $p$ is odd, $v_p(w)<0$, and both local slots of $W^{-1}$ are
integral. Then $\mathcal R_a\equiv1\pmod p$, provided $a$ is a
$p$-adic unit. In particular this applies at an ordinary conditional
pole $t\to1$ whenever the limiting product
\[
                     \prod_{i=1}^6(1+C_i^2/7)
\]
is nonzero and a local unit in both slots.
\end{lemma}
\begin{proof}
The trace and norm $V,J$ are integral. Since $v_p(g)=v_p(w)<0$,
one has $Q\equiv1$, and $1+a^2(1-V/2)/w\equiv1\pmod p$.
At the specified conditional pole $F(1)=7$ and the displayed unit
condition gives integrality of the reciprocal slots.
\end{proof}

These are local statements for the stated actual preparation patterns.
They do not assert local square values at every prime. In particular,
new prime divisors of detector numerators at even physical valuations
remain subject to the square condition itself. Nor does local
solubility imply simultaneous rational square values.

\section{Finite algebra and the remaining global realization}

\begin{definition}\label{def:realization}
The simultaneous native square property for a source means that every
required integer input admits a prepared actual rational source tuple
on the whole-character real and guard locus specified above, with all
inherited point, difference, scalar, trace and comparison guards and
the unchanged common slope and height scales, such that
\[
                       \mathcal R_{a_0,0},\quad
                       \mathcal R_{a_1,1}\in\Q^{\times2}.
\]
For a conditional source the required inputs are the nonzero integers.
Each $a_j$ is fixed in advance either as a nonzero $S$-unit or as an
admissible rational-function expression in the existing source
coordinates, as in Proposition~\ref{prop:adaptive} or
Corollary~\ref{cor:outer}; every denominator of that expression is
included in the retained guard locus. It is not a new unrestricted
parameter. These subscripts are independent of the source's notation
for its comparison scales.
\end{definition}

\begin{theorem}[Guarded conditional reduction]\label{thm:count}
Assume the retained arithmetic premises of the normalized general
source or one of the conditional sources, and assume its simultaneous
native square property. Then the corresponding defining polynomial
may be constructed with twenty-one total variables in the general
case, and nine in either conditional case, including the input.

Without that realization assumption, the square criterion proves the
soundness direction of this finite-algebra construction only. It does
not improve the established defining counts.
\end{theorem}
\begin{proof}
Retain $p$ point pairs, with $p=9$ in the general source and $p=3$
in either conditional source. Their finite point algebra has rank
$4^p$, as in~\cite{Unit}. Replace each pole block by its square test.
Write the two actual radicands as $C_j/D_j$ in the localized point
algebra and adjoin the two monic finite generators
\[
                     Y_j^2=C_jD_j,\qquad j=0,1.
\]
Then the square root is $Y_j/D_j$. At a prepared tuple the numerator
and denominator norms are nonzero by Theorem~\ref{thm:criterion}
and the inherited guard audit. Include these scalar norms in the guard
product. The two quadratic extensions are finite etale over the guarded
point algebra, have radicands in that original algebra, and give total rank
\[
                              4^p\,2^2=4^{p+1}.
\]
The one-coordinate guarded separator and determinant descent of
\cite{Unit} therefore apply without an additional unrestricted
parameter. Soundness follows from Theorem~\ref{thm:criterion} and the
unchanged source local extraction and height argument. The realization
assumption extends the selected rational point character to both
quadratic generators, proving completeness.

The unrestricted coordinates are the input, $2p$ point coordinates,
one common slope and one final determinant coordinate, totaling
$2p+3$. In a conditional source the usual final input factor includes
zero. This gives the stated counts.
\end{proof}

Definition~\ref{def:realization} is the unresolved arithmetic step.
The local results above preserve prescribed actual traces and can
filter candidate families, but neither their congruences nor the
reference identity~\eqref{eq:baseline} establish that definition.
An affirmative result must produce rational points on both actual
square covers at one common slope, with all point lifts and guards.

\begin{thebibliography}{9}
\raggedright
\bibitem{General}
\emph{A defining polynomial in twenty-five variables: combined scalar
and comparison conditions}, repository source,
\nolinkurl{papers/few_variables25/few_variables25_merged_20261110_v19.tex}.
\bibitem{Conditional}
\emph{Thirteen variables from direct elliptic index comparison},
repository source,
\nolinkurl{papers/conditional_few_variables/direct_input13_20261110_v22.tex}.
\bibitem{Unit}
\emph{A single norm determinant for the 25- and 13-variable constructions},
repository source,
\nolinkurl{papers/finite_algebra/unit_norm_20261114_v28.tex}.
\bibitem{Normalized}
\emph{Quadratic reweightings of native elliptic detectors},
companion source specification, 16 November 2026,
\nolinkurl{papers/pole_parameters/reweighted_elliptic_detectors_20261116_v53.tex}.
\end{thebibliography}
\end{document}
